\documentclass{Config/EngineeringThesis}   % options: auxsup, secsup, pl, en
\usepackage{enumerate}
\usepackage{enumitem}

\usepackage{caption}
\usepackage{indentfirst}
\usepackage[a4paper, left=25mm, top=25mm, right=25mm, bottom=25mm]{geometry} % marginesy 2,5 cm (wytyczne redakcyjne AGH)

\usepackage{amsmath,amsthm,amssymb,amsfonts}
\usepackage{titlesec}
\usepackage[polish]{babel} 
\addto\captionsenglish{\renewcommand{\figurename}{Fig.}}
\addto\captionspolish{\renewcommand{\figurename}{Rys.}}
\usepackage[labelsep=period]{caption}
\usepackage[format=plain, labelfont={it}, textfont=it]{caption}
\usepackage{notoccite}
\usepackage{tocloft}
\setlength{\cftbeforetoctitleskip}{-3em}
\usepackage[inkscapelatex=false]{svg}

% Tables
\usepackage{array}
\usepackage{longtable}
\newcolumntype{P}[1]{>{\centering\arraybackslash}p{#1}}
\newcolumntype{M}[1]{>{\centering\arraybackslash}m{#1}}
\newcolumntype{C}[1]{>{\centering\arraybackslash}p{#1}}
\newcolumntype{L}[1]{>{\raggedright\arraybackslash}p{#1}}
\usepackage{makecell}
\usepackage{diagbox}
\usepackage{caption, threeparttable}

\usepackage{subcaption}
\usepackage{graphicx}           
\usepackage{lipsum} 
\usepackage{environ} 
\usepackage{float}
\usepackage{pdfpages}


\usepackage{tikz}
\usetikzlibrary{arrows}
\usetikzlibrary{arrows.meta}
\usepackage{circuitikz}               % schematy elektryczne
\usepackage{pgfplots}                 % wykresy
\usepgfplotslibrary{groupplots}
\pgfplotsset{compat=1.17}
\usepackage{tabularray}
\usepackage{siunitx,booktabs,array}

\usepackage{amsmath}
\usepackage{arydshln}
\usepackage{bigints}
\usepackage{relsize}
\usepackage{mathtools}
\usepackage{moresize}
\usepackage{nicefrac}

\usepackage{pdflscape}
\usepackage{keyval}

\usepackage{textcomp}
\usepackage{siunitx}
\sisetup{locale = PL}

\usepackage[numbers,sort&compress]{natbib}

% Linia ze źródłem pod rysunkiem/tabelą: \zrodlo{opracowanie własne}
\newcommand{\zrodlo}[1]{%
  \par\vspace{3pt}{\centering\footnotesize Źródło: #1\par}}

\def  \EQNSCALE {1.1}
\def  \EQNTAB {1.1}
\NewEnviron{EQNSIZE}{% 
    \scalebox{\EQNSCALE}{$\BODY$} 
}

%%%%%%%%%%%%%%%%%%%%%%%%%%%%%%%--------------------------------xxxxxxxxxxxxxxxxxxxxxxxxxxxxxxxx






% ###################### Chapter name formatting ##############################################

\makeatletter
\def\@makechapterhead#1{
  {\parindent \z@ \raggedright \normalfont
    \ifnum \c@secnumdepth >\m@ne
        %\Huge\bfseries \MakeUppercase\@chapapp\space \thechapter:   \space
        \Huge\bfseries \MakeUppercase \thechapter.   \space
        %\par\nobreak
        %\vskip 20\p@
    \fi
    \interlinepenalty\@M
    \huge \bfseries #1\par\nobreak
    %\vskip 20\p@
  }}
\makeatother

% ###################### End of chapter name formatting ########################################

% polish if the pl option is used
\usepackage[nottoc,notlof,notlot]{tocbibind}
%\usepackage[nottoc]{tocbibind}
\usepackage{multirow}
\usepackage[hidelinks]{hyperref}
\hypersetup{
    colorlinks=false,
    linkcolor=black,
    filecolor=black,      
    urlcolor=black
}

\theoremstyle{definition}
\newtheorem{theorem}{Theorem}[chapter]
\newtheorem{proposition}[theorem]{Proposition}
\newtheorem{lemma}[theorem]{Lemma}
\newtheorem{definition}{Definition}[chapter]

\graphicspath{{./Pictures/}}




%---------------------------------------------------------------------------
% University name:
\universitynamepl {Akademia Górniczo-Hutnicza im. Stanisława Staszica w Krakowie}
\universitynameen {AGH University of Krakow}

% Faculty name:
\facultynamepl {Wydział Elektrotechniki, Automatyki, Informatyki i Inżynierii Biomedycznej}
\facultynameen {Faculty of Electrical Engineering, Automatics, Computer Science and Biomedical Engineering}

% field of study:
\fieldofstudypl {Elektrotechnika}
\fieldofstudyen {Electrical Engineering}

% title of the engineering project:
\titlepl{Tytuł projektu inżynierskiego}
\titleen{Project title in English}
 


% Author names and surname: 
\authorpl{Jan Antoni Nowak}
\authoren{Jan Antoni Nowak}

% first supervisor name with degrees:
\supervisorpl{dr inż. Jan Nowak}
\supervisoren{Jan Nowak, PhD}

% second supervisor (or auxiliary supervisor) name with degress:
%\secondsupervisor{Prof. dr hab. inż. Adam Iksiński} 
\secondsupervisorpl{}
\secondsupervisoren{}

\date{202x}

\acknowledgementspl{Specjalne podziękowania dla \dots}
\acknowledgementsen{Special thanks to \dots}

\setlength{\cftsecnumwidth}{10mm}
\allowdisplaybreaks

\usepackage{listings}
\usepackage{xcolor}

\definecolor{codegreen}{rgb}{0,0.6,0}
\definecolor{codegray}{rgb}{0.5,0.5,0.5}
\definecolor{codepurple}{rgb}{0.58,0,0.82}
\definecolor{backcolour}{rgb}{1,1,1}

\lstdefinestyle{mystyle_listing}{
    backgroundcolor=\color{backcolour},   
    commentstyle=\color{codegreen},
    keywordstyle=\color{magenta},
    numberstyle=\tiny\color{codegray},
    stringstyle=\color{codepurple},
    basicstyle=\ttfamily\footnotesize,
    breakatwhitespace=false,         
    breaklines=true,                 
    captionpos=b,                    
    keepspaces=true,                 
    numbers=left,                    
    numbersep=5pt,                  
    showspaces=false,                
    showstringspaces=false,
    showtabs=false,                  
    tabsize=2,
    frame = single
}

\lstset{style=mystyle_listing}

\lstset{
  literate={ą}{{\k a}}1
  		     {Ą}{{\k A}}1
           {ż}{{\. z}}1
           {Ż}{{\. Z}}1
           {ź}{{\' z}}1
           {Ź}{{\' Z}}1
           {ć}{{\' c}}1
           {Ć}{{\' C}}1
           {ę}{{\k e}}1
           {Ę}{{\k E}}1
           {ó}{{\' o}}1
           {Ó}{{\' O}}1
           {ń}{{\' n}}1
           {Ń}{{\' N}}1
           {ś}{{\' s}}1
           {Ś}{{\' S}}1
           {ł}{{\l}}1
           {Ł}{{\L}}1
}

% Lokalizacja w źródle w \cite[s.~200]{klucz} służy tylko autorowi do
% weryfikacji i nie jest drukowana w PDF. Aby ją wydrukować, zamień
% \pokazstronyfalse na \pokazstronytrue.
\newif\ifpokazstrony
\pokazstronyfalse
\AtBeginDocument{%
  \ifpokazstrony\else
    \let\citezlokalizacja\cite
    \renewcommand{\cite}[2][]{\citezlokalizacja{#2}}%
  \fi}

\begin{document}

\frontmatter

\thesistitlepagepl
%\thesistitlepageen

\pagenumbering{gobble}
\include{abstract-pl}  % abstract in English 
\pagenumbering{gobble}
\include{abstract-en}  % abstract in Polish
\cleardoublepage

\makeatletter
\renewcommand\tableofcontents{%
    \if@twocolumn
      \@restonecoltrue\onecolumn
    \else
      \@restonecolfalse
    \fi
    \chapter*{\vspace{0cm}\contentsname
        \@mkboth{%
           \MakeUppercase\contentsname}{\MakeUppercase\contentsname}}%
    \@starttoc{toc}%
    \if@restonecol\twocolumn\fi
    }
\makeatother

\tableofcontents%\thispagestyle{empty}\newpage

\cleardoublepage
%\listoffigures
%\cleardoublepage

%**********************************************************
\renewcommand{\chaptermark}[1]{\markboth{\thechapter:\ #1}{}}
\renewcommand{\sectionmark}[1]{\markright{\thesection.\ #1}}
\renewcommand{\headrulewidth}{0.5pt}
\renewcommand{\footrulewidth}{0.5pt}
\fancyhf{}
\fancyhead{}
\fancyhead[OR]{\small \MakeUppercase{\chaptername}\space\leftmark}
\fancyhead[EL]{\small \MakeUppercase{\chaptername}\space\leftmark}
%\fancyhead[L]{\small \chaptername \space \thechapter: \chaptertitlename}
%\fancyhead[R]{\bfseries \thepage}
\fancyfoot[C]{\bfseries \thepage}
%\thesisfoot

\newenvironment{conditions}
{\par\vspace{5pt}\begin{tabular}{>{$}l<{$} @{${}-{}$} l}}
  {\end{tabular}\par\vspace{\belowdisplayskip}}
  
  %{\par\vspace{\abovedisplayskip}\noindent\begin{tabular}{>{$}l<{$} @{${}-{}$} l}}
  %{\end{tabular}\par\vspace{\belowdisplayskip}}



%**************************************************************

\renewcommand{\UrlFont}{\normalfont}


\mainmatter
\addtocontents{toc}{\protect\setcounter{tocdepth}{2}}


% add your chapter here
\begin{titlepage}
%\pagestyle{plain}
\pagenumbering{gobble}
\pagestyle{empty}
\titleformat{\chapter}[display]{\normalfont\bfseries}{}{0pt}{\Huge}
\titlespacing*{\chapter}{0pt}{2cm}{0pt}
\chapter*{Wykaz akronimów}
\addcontentsline{toc}{chapter}{Wykaz akronimów}
\setcounter{chapter}{0}

\begin{table}[H]
\begin{longtable}{ C{30mm}  C{10mm}   C{107mm} }
\multicolumn{1}{c}{\makecell[tl]{DC-DC}} & \multicolumn{1}{c}{\makecell[tc]{---}} & \multicolumn{1}{c}{\makecell[tl]{ang. Direct Current to Direct Current Converter\\ (przekształtnik prądu stałego na stały)\\}} \\
\multicolumn{1}{c}{\makecell[tl]{LED}} & \multicolumn{1}{c}{\makecell[tc]{---}} & \multicolumn{1}{c}{\makecell[tl]{ang. Light Emitting Diode \\ (dioda elektroluminescencyjna)\\}} \\
\multicolumn{1}{c}{\makecell[tl]{CPU}} & \multicolumn{1}{c}{\makecell[tc]{---}} & \multicolumn{1}{c}{\makecell[tl]{ang. Central Processing Unit \\ (procesor)\\}} \\
\multicolumn{1}{c}{\makecell[tl]{GPU}} & \multicolumn{1}{c}{\makecell[tc]{---}} & \multicolumn{1}{c}{\makecell[tl]{ang. Graphics Processing Unit \\ (procesor graficzny)\\}} \\
\multicolumn{1}{c}{\makecell[tl]{PV}} & \multicolumn{1}{c}{\makecell[tc]{---}} & \multicolumn{1}{c}{\makecell[tl]{ang. Photovoltaics \\ (fotowoltaika)\\}} \\
\multicolumn{1}{c}{\makecell[tl]{MOSFET}} & \multicolumn{1}{c}{\makecell[tc]{---}} & \multicolumn{1}{c}{\makecell[tl]{ang. Metal-Oxide-Semiconductor Field-Effect \\ Transistor (tranzystor polowy z izolowaną bramką)\\}} \\
\multicolumn{1}{c}{\makecell[tl]{GaN}} & \multicolumn{1}{c}{\makecell[tc]{---}} & \multicolumn{1}{c}{\makecell[tl]{azotek galu - nieorganiczny związek chemiczny \\ (półprzewodnik)\\}} \\
\multicolumn{1}{c}{\makecell[tl]{CoolMOS}} & \multicolumn{1}{c}{\makecell[tc]{---}} & \multicolumn{1}{c}{\makecell[tl]{Rodzina tranzystorów krzemowych (Si) \\ produkowanych przez firmę Infineon.\\}} \\
\multicolumn{1}{c}{\makecell[tl]{CoolGaN}} & \multicolumn{1}{c}{\makecell[tc]{---}} & \multicolumn{1}{c}{\makecell[tl]{Rodzina tranzystorów z azotku galu (GaN) \\ produkowanych przez firmę Infineon.\\}} \\
\multicolumn{1}{c}{\makecell[tl]{IoT}} & \multicolumn{1}{c}{\makecell[tc]{---}} & \multicolumn{1}{c}{\makecell[tl]{ang. Internet of Things \\ (Sieć połączonych ze sobą urządzeń zdolnych\\ do współdzielenia i wymiany informacji.)\\}} \\
\multicolumn{1}{c}{\makecell[tl]{LCD}} & \multicolumn{1}{c}{\makecell[tc]{---}} & \multicolumn{1}{c}{\makecell[tl]{ang. Liquid Crystal Display \\ (wyświetlacz ciekłokrystaliczny)\\}} \\
\multicolumn{1}{c}{\makecell[tl]{HVGDCC}} & \multicolumn{1}{c}{\makecell[tc]{---}} & \multicolumn{1}{c}{\makecell[tl]{ang. High Voltage Gain Direct Current Converter\\ (Topologia wielopoziomowego przekształtnika\\ DC-DC o przełącznych kondensatorach, opracowana \\ w ramach prac związanych z przygotowaniem \\ niniejszej rozprawy doktorskiej.)\\}} \\

 &  &    \\ % This line must stay empty. Their lack cause issues with table aligment. 
\end{longtable}
\end{table}

\begin{table}[H]
\begin{longtable}{ C{30mm}  C{10mm}   C{107mm} }
\multicolumn{1}{c}{\makecell[tl]{HVGDCC-CoolMOS}} & \multicolumn{1}{c}{\makecell[tc]{---}} & \multicolumn{1}{c}{\makecell[tl]{Przekształtnik o topologii HVGDCC, którego obwód \\mocy zbudowano z wykorzystaniem tranzystorów \\ typu CoolMOS\\}} \\
\multicolumn{1}{c}{\makecell[tl]{HVGDCC-CoolGaN}} & \multicolumn{1}{c}{\makecell[tc]{---}} & \multicolumn{1}{c}{\makecell[tl]{Przekształtnik o topologii HVGDCC, którego obwód \\mocy zbudowano z wykorzystaniem tranzystorów \\ typu CoolGaN\\}} \\
\multicolumn{1}{c}{\makecell[tl]{BMS}} & \multicolumn{1}{c}{\makecell[tc]{---}} & \multicolumn{1}{c}{\makecell[tl]{ang. Battery Management System \\(Układ wyrównujący napięcia na poszczególnych \\ ogniwach akumulatora)\\}} \\
 &  &    \\ % This line must stay empty. Their lack cause issues with table aligment. 
\end{longtable}
\end{table}

\end{titlepage}
\newpage
\pagenumbering{arabic}
\pagestyle{plain}
\titleformat{\chapter}[display]{\normalfont\bfseries}{}{0pt}{\Huge}
\titlespacing*{\chapter}{0pt}{2cm}{0pt}
\chapter*{Wykaz oznaczeń}
\addcontentsline{toc}{chapter}{Wykaz oznaczeń}
\setcounter{chapter}{0}

\begin{table}[H]
\begin{longtable}{ C{30mm}  C{10mm}   C{107mm} }
\multicolumn{1}{c}{\makecell[tl]{HVGDCC}} & \multicolumn{1}{c}{\makecell[tc]{---}} & \multicolumn{1}{c}{\makecell[tl]{ang. High Voltage Gain Direc Current Converter \\ dsddsvsd\\ \\ }} \\
\multicolumn{1}{c}{\makecell[tl]{HVGDCC}} & \multicolumn{1}{c}{\makecell[tc]{---}} & \multicolumn{1}{c}{\makecell[tl]{ang. High Voltage Gain Direc Current Converter \\ dsddsvsd\\ \\}} \\
\multicolumn{1}{c}{\makecell[tl]{HVGDCC}} & \multicolumn{1}{c}{\makecell[tc]{---}} & \multicolumn{1}{c}{\makecell[tl]{ang. High Voltage Gain Direc Current Converter \\ dsddsvsd\\ \\}} \\
 &  &    \\ % This line must stay empty. Their lack cause issues with table aligment. 
\end{longtable}
\end{table}
\chapter{Wprowadzenie}
\label{chap: Wprowadzenie}


\section{Cel i zakres pracy}
\label{sec: Cel i zakres pracy}

\section{Założenia projektowe}
\label{sec: Założenia projektowe}

\section{Struktura pracy}
\label{sec: Struktura pracy}
\chapter{Falowniki jednofazowe - podstawy teoretyczne}
\label{chap: Falowniki jednofazowe - podstawy teoretyczne}

\section{Zasada działania falowników napięcia}
\label{sec: Zasada działania falowników napięcia}

\section{Topologie jednofazowych falowników}
\label{sec: Topologie jednofazowych falowników}

\section{Metody modulacji PWM}
\label{sec: Metody modulacji PWM}

\section{Zastosowania falowników w energoelektronice}
\label{sec: Zastosowania falowników w energoelektronice}
\chapter{Charakterystyka projektowanego stanowiska laboratoryjnego}
\label{chap: Charakterystyka projektowanego stanowiska laboratoryjnego}

W~rozdziale scharakteryzowano projektowane stanowisko laboratoryjne
z~jednofazowym falownikiem napięcia. Określono wymagania funkcjonalne,
parametry elektryczne układu, wymagania dotyczące bezpieczeństwa
użytkowania oraz wymagania dydaktyczne, a~następnie przedstawiono
architekturę stanowiska, wynikającą z~tych wymagań. Sformułowane tu
założenia stanowią punkt wyjścia do doboru elementów obwodu mocy, opisanego
w~rozdziale~\ref{chap: Dobór elementów obwodu mocy falownika}.

% ---------------------------------------------------------------------
\section{Wymagania funkcjonalne}
\label{sec: Wymagania funkcjonalne}

Projektowane stanowisko ma umożliwić praktyczne zapoznanie się z~pracą
jednofazowego falownika napięcia w~układzie mostka~H oraz doświadczalną
weryfikację zjawisk omówionych w~rozdziale~\ref{chap: Falowniki jednofazowe - podstawy teoretyczne}.
Stanowi ono uzupełnienie zajęć symulacyjnych z~falowników
jednofazowych~\cite{instrukcja}, w~trakcie których badane są modele
falowników sterowanych metodą przebiegu prostokątnego, metodą modulacji
szerokości pojedynczego impulsu w~półokresie oraz metodami bipolarnej
i~unipolarnej modulacji SPWM. Stanowisko powinno pozwalać na
odtworzenie tych samych przebiegów w~rzeczywistym układzie, a~także na
obserwację zjawisk pomijanych w~modelach idealnych, takich jak procesy
łączeniowe tranzystorów, wpływ czasu martwego, straty mocy czy zakłócenia
związane z~szybkimi zmianami napięcia.

W~wymaganiach funkcjonalnych określono, \emph{jakie funkcje} ma realizować
stanowisko, bez wskazywania konkretnych wartości liczbowych. Parametry
elektryczne zestawiono w~podrozdziale~\ref{sec: Parametry elektryczne projektowanego układu},
a~szczegółowe wymagania dotyczące bezpieczeństwa i~dydaktyki --
w~podrozdziałach~\ref{sec: Wymagania dotyczące bezpieczeństwa użytkowania}
i~\ref{sec: Wymagania dydaktyczne}. Przyjęto następujące wymagania
funkcjonalne.

\begin{enumerate}[label=\textbf{F\arabic*.}, leftmargin=*, itemsep=4pt]

  \item \textbf{Przekształcanie energii prądu stałego na energię prądu
    przemiennego.} Stanowisko ma pracować jako falownik napięcia w~układzie
    pełnego mostka (dwie gałęzie łącznikowe, tranzystory $Q_1$--$Q_4$).
    Obwód pośredniczący należy zasilać z~zewnętrznego zasilacza
    laboratoryjnego o~regulowanym napięciu, a~nie bezpośrednio z~sieci
    elektroenergetycznej. Pozwala to rozpoczynać pracę od niskiego napięcia
    $U_\mathrm{DC}$ i~zwiększać je stopniowo oraz ograniczać prąd pobierany
    przez układ nastawą zasilacza. Ponieważ zasilacz laboratoryjny zwykle
    nie może przyjmować energii, energię zwracaną z~obciążenia
    o~charakterze indukcyjnym musi przejmować kondensator obwodu
    pośredniczącego $C_\mathrm{DC}$.

  \item \textbf{Niezależne sterowanie każdym z~czterech tranzystorów.}
    Każdy tranzystor mostka ma mieć własny, galwanicznie odizolowany kanał
    sterowania bramkowego, zasilany z~osobnego izolowanego źródła. Takie
    rozwiązanie, w~odróżnieniu od zasilania typu \textit{bootstrap}, nie
    nakłada ograniczeń na współczynnik wypełnienia ani na czas, przez jaki
    tranzystor górny może pozostawać załączony. Umożliwia to realizację
    dowolnego schematu łączeń, w~tym sterowania przebiegiem prostokątnym
    o~częstotliwości $f_1$, w~którym każdy tranzystor przewodzi przez
    połowę okresu napięcia wyjściowego, a~także pracę statyczną, np.~przy
    sprawdzaniu poszczególnych tranzystorów.

  \item \textbf{Realizacja różnych metod sterowania.} Stanowisko ma
    umożliwiać sterowanie falownikiem co najmniej metodami opisanymi
    w~podrozdziale~\ref{sec: Metody modulacji PWM}: przebiegiem
    prostokątnym, modulacją szerokości pojedynczego impulsu (przez
    przesunięcie fazowe sygnałów sterujących gałęzi) oraz bipolarną
    i~unipolarną modulacją SPWM. Zmiana metody sterowania nie powinna
    wymagać zmian w~obwodzie mocy, a~jedynie zmiany sygnałów sterujących.

  \item \textbf{Regulacja parametrów sterowania.} Użytkownik ma mieć
    możliwość zmiany częstotliwości podstawowej $f_1$, współczynnika
    modulacji amplitudowej $m_\mathrm{a}$, częstotliwości łączeń
    $f_\mathrm{s}$ oraz czasu martwego $t_\mathrm{d}$. Regulacja
    $m_\mathrm{a}$ powinna obejmować również zakres nadmodulacji, a~regulacja
    $t_\mathrm{d}$ -- wartości wyraźnie większe od minimalnych, tak aby
    można było obserwować zniekształcenia napięcia wyjściowego wywołane
    czasem martwym.

  \item \textbf{Oddzielenie części sterującej od obwodu mocy.} Sygnały
    sterujące tranzystorami mają być doprowadzane do stanowiska jako
    sygnały logiczne z~zewnętrznego układu sterującego (mikrokontrolera lub
    generatora sygnałów). Wejścia sterujące oraz obwody pomiarowe
    i~zabezpieczające po stronie układu sterującego mają być galwanicznie
    odizolowane od obwodu mocy. Dzięki temu układ sterujący, komputer
    i~przyrządy pomiarowe użytkownika nie są narażone na potencjały
    występujące w~mostku.

  \item \textbf{Zasilanie różnych odbiorników.} Do wyjścia falownika ma
    być możliwe dołączenie obciążenia rezystancyjnego,
    rezystancyjno-indukcyjnego oraz filtru wyjściowego LC. Pozwala to
    porównać przebiegi prądu wyjściowego $i_\mathrm{o}$ i~napięcia
    wyjściowego $u_\mathrm{o}$ przed filtrem i~za nim oraz zbadać wpływ
    indukcyjności $L_\mathrm{f}$ na tętnienia prądu.

  \item \textbf{Dostęp do sygnałów pomiarowych.} Stanowisko ma zapewniać
    łatwy dostęp do punktów pomiarowych, umożliwiających rejestrację
    oscyloskopem co najmniej: napięć gałęzi $u_\mathrm{aN}$
    i~$u_\mathrm{bN}$, napięcia wyjściowego $u_\mathrm{o}$, prądu
    wyjściowego $i_\mathrm{o}$, napięcia obwodu pośredniczącego
    $U_\mathrm{DC}$ oraz sygnałów sterujących tranzystorami. Prąd
    wyjściowy ma być mierzony w~układzie wbudowanym, dającym sygnał
    napięciowy proporcjonalny do prądu, co pozwala obserwować go bez
    sondy prądowej.

  \item \textbf{Możliwość pomiaru sprawności.} Zaciski wejściowe
    i~wyjściowe falownika mają umożliwiać dołączenie przyrządów do pomiaru
    mocy pobieranej z~zasilacza i~mocy oddawanej do obciążenia, tak aby
    można było wyznaczyć sprawność przekształtnika w~funkcji obciążenia.

  \item \textbf{Zabezpieczenia i~sygnalizacja stanu.} Stanowisko ma
    samoczynnie blokować pracę tranzystorów po przekroczeniu
    dopuszczalnego prądu, niezależnie od oprogramowania układu
    sterującego, oraz informować układ sterujący o~zadziałaniu
    zabezpieczenia. Tranzystory nie mogą być załączane, gdy napięcie
    zasilania sterowników bramkowych jest zbyt niskie. Po wyłączeniu
    stanowiska musi być możliwe szybkie rozładowanie kondensatora obwodu
    pośredniczącego przez rezystory rozładowcze, załączane przyciskiem,
    a~stan pracy układu (obecność napięcia, zadziałanie zabezpieczeń) ma
    być sygnalizowany optycznie. Szczegółowe wymagania
    przedstawiono w~podrozdziale~\ref{sec: Wymagania dotyczące bezpieczeństwa użytkowania}.

  \item \textbf{Łatwa wymiana łączników.} Tranzystory mocy mają być
    montowane w~sposób umożliwiający ich wymianę bez lutowania. Pozwala to
    szybko wymienić uszkodzony element, a~także porównywać pracę
    falownika z~tranzystorami różnych typów.

\end{enumerate}

Zestawienie wymagań funkcjonalnych wraz ze wskazaniem miejsc pracy,
w~których opisano sposób ich realizacji i~weryfikacji, przedstawiono
w~tabeli~\ref{tab:wymagania-funkcjonalne}.

\begin{table}[htb]
  \centering
  \caption{Zestawienie wymagań funkcjonalnych projektowanego stanowiska}
  \label{tab:wymagania-funkcjonalne}
  \small
  \begin{tabular}{@{}C{10mm} L{80mm} L{55mm}@{}}
    \toprule
    \textbf{Nr} & \textbf{Wymaganie} & \textbf{Realizacja / weryfikacja} \\
    \midrule
    F1  & Przekształcanie DC/AC w~układzie mostka~H, zasilanie z~zasilacza
          laboratoryjnego
        & podrozdz.~\ref{sec: Architektura projektowanego stanowiska},
          \ref{sec: Dobór kondensatorów magistrali DC} \\
    F2  & Niezależne, izolowane sterowanie każdego z~tranzystorów $Q_1$--$Q_4$
        & podrozdz.~\ref{sec: Dobór sterowników bramkowych} \\
    F3  & Sterowanie przebiegiem prostokątnym, pojedynczym impulsem oraz
          SPWM bipolarną i~unipolarną
        & podrozdz.~\ref{sec: Analiza przebiegów napięć i prądów},
          \ref{sec: Przebiegi napięć i prądów} \\
    F4  & Regulacja $f_1$, $m_\mathrm{a}$, $f_\mathrm{s}$ i~$t_\mathrm{d}$
        & podrozdz.~\ref{sec: Możliwe ćwiczenia laboratoryjne} \\
    F5  & Izolacja galwaniczna części sterującej od obwodu mocy
        & podrozdz.~\ref{sec: Architektura projektowanego stanowiska},
          \ref{sec: Projekt obwodu drukowanego} \\
    F6  & Obciążenie R, RL oraz filtr wyjściowy LC
        & podrozdz.~\ref{sec: Przebiegi napięć i prądów} \\
    F7  & Punkty pomiarowe napięć, prądu i~sygnałów sterujących
        & podrozdz.~\ref{sec: Projekt obwodu drukowanego},
          \ref{sec: Stanowisko pomiarowe} \\
    F8  & Pomiar sprawności przekształtnika
        & podrozdz.~\ref{sec: Badanie sprawności przekształtnika} \\
    F9  & Zabezpieczenie nadprądowe i~podnapięciowe, rozładowanie
          $C_\mathrm{DC}$, sygnalizacja stanu
        & podrozdz.~\ref{sec: Wymagania dotyczące bezpieczeństwa użytkowania},
          \ref{sec: Ocena działania zabezpieczeń} \\
    F10 & Wymiana tranzystorów bez lutowania
        & podrozdz.~\ref{sec: Projekt obwodu drukowanego} \\
    \bottomrule
  \end{tabular}
  \zrodlo{opracowanie własne}
\end{table}


\section{Parametry elektryczne projektowanego układu}
\label{sec: Parametry elektryczne projektowanego układu}

\section{Wymagania dotyczące bezpieczeństwa użytkowania}
\label{sec: Wymagania dotyczące bezpieczeństwa użytkowania}

\section{Wymagania dydaktyczne}
\label{sec: Wymagania dydaktyczne}

\section{Architektura projektowanego stanowiska}
\label{sec: Architektura projektowanego stanowiska}

\chapter{Dobór elementów obwodu mocy falownika}
\label{chap: Dobór elementów obwodu mocy falownika}

\section{Dobór tranzystorów MOSFET}
\label{sec: Dobór tranzystorów MOSFET}

\section{Dobór sterowników bramkowych}
\label{sec: Dobór sterowników bramkowych}

Jak wskazano w~podrozdziale~\ref{sec: Topologie jednofazowych falowników},
tranzystory górne mostka~H wymagają sterowników o~pływającym potencjale
odniesienia. Potencjał węzłów wyjściowych zmienia się przy każdym przełączeniu
między potencjałami biegunów N~i~P z~szybkością narastania rzędu kilku do
kilkudziesięciu woltów na nanosekundę. Aby załączyć tranzystor górny,
napięcie jego bramki musi być o~$U_\mathrm{GS(on)}$ wyższe od potencjału węzła
wyjściowego, czyli względem bieguna N wynosić nawet
$U_\mathrm{DC} + U_\mathrm{GS(on)}$. Wynikają z~tego trzy wymagania wobec
toru sterowania bramkowego:
\begin{itemize}
  \item \textbf{Zasilanie odniesione do źródła tranzystora.} Napięcie
    zasilające sterownik musi być podane względem źródła tranzystora górnego,
    a~nie względem masy układu~\cite[s.~29, 142]{mohan2012}. Stosuje się
    w~tym celu układ bootstrap albo izolowaną przetwornicę DC/DC
    z~transformatorem~\cite[s.~15]{instrukcja}, \cite[rozdz.~5, 7]{balogh2018}.
  \item \textbf{Przeniesienie sygnału sterującego.} Sygnał PWM z~mikrokontrolera,
    odniesiony do masy układu sterowania, trzeba przenieść do sterownika
    o~pływającym potencjale -- przez układ przesuwania poziomu
    (ang.~\textit{level shifter}) albo przez izolację galwaniczną (izolator
    cyfrowy, transoptor)~\cite[s.~29]{mohan2012}, \cite[rozdz.~5]{balogh2018}.
  \item \textbf{Odporność na szybkie zmiany potencjału.} Pojemności
    pasożytnicze między stroną pływającą a~masą układu sterowania (pojemność
    międzyuzwojeniowa transformatora przetwornicy, pojemność bariery
    izolatora) przewodzą przy każdym przełączeniu prąd
    $i = C\,\mathrm{d}u/\mathrm{d}t$. Przy pojemności \SI{10}{\pico\farad}
    i~szybkości narastania \SI{50}{\volt\per\nano\second} jest to impuls
    prądu o~wartości \SI{0,5}{\ampere}, który może zakłócić pracę sterownika
    lub mikrokontrolera. Elementy izolujące powinny więc mieć możliwie małą
    pojemność sprzęgającą.
\end{itemize}

Układ bootstrap jest prosty i~tani: kondensator zasilający sterownik górny
ładuje się przez diodę z~zasilania strony dolnej w~czasie, gdy przewodzi
tranzystor dolny tej samej gałęzi. Z~tego powodu nie pozwala on na długotrwałe
załączenie tranzystora górnego (współczynnik wypełnienia bliski 1), a~uzyskanie
ujemnego napięcia wyłączającego wymaga dodatkowych elementów~\cite[rozdz.~5]{balogh2018}.
Izolowane zasilanie każdego kanału jest droższe i~zajmuje więcej miejsca, ale
umożliwia pracę z~dowolnym współczynnikiem wypełnienia (również stałe
załączenie łącznika) oraz stosowanie ujemnego napięcia wyłączającego bramkę,
co zwiększa odporność na pasożytnicze załączenie tranzystora. Z~tych powodów
w~projektowanym stanowisku zastosowano izolowane zasilanie każdego z~czterech
kanałów sterowania bramkowego.


\section{Dobór kondensatorów magistrali DC}
\label{sec: Dobór kondensatorów magistrali DC}

\section{Analiza strat mocy}
\label{sec: Analiza strat mocy}

\section{Dobór radiatora}
\label{sec: Dobór radiatora}

\chapter{Symulacje}
\label{chap: Symulacje}

\section{Opracowanie modelu symulacyjnego}
\label{sec: Opracowanie modelu symulacyjnego}

\section{Analiza przebiegów napięć i prądów}
\label{sec: Analiza przebiegów napięć i prądów}

\section{Weryfikacja poprawności doboru komponentów falownika}
\label{sec: Weryfikacja poprawności doboru komponentów falownika}

\chapter{Wykonanie modelu falownika}
\label{chap: Wykonanie modelu falownika}

\section{Projekt obwodu drukowanego}
\label{sec: Projekt obwodu drukowanego}

\section{Uruchomienie układu}
\label{sec: Uruchomienie układu}

\section{Napotkane problemy konstrukcyjne i wprowadzone modyfikacje}
\label{sec: Napotkane problemy konstrukcyjne i wprowadzone modyfikacje}
\chapter{Badania eksperymentalne}
\label{chap: Badania eksperymentalne}

\section{Stanowisko pomiarowe}
\label{sec: Stanowisko pomiarowe}

\section{Przebiegi napięć i prądów}
\label{sec: Przebiegi napięć i prądów}

\section{Badanie sprawności przekształtnika}
\label{sec: Badanie sprawności przekształtnika}

\section{Ocena działania zabezpieczeń}
\label{sec: Ocena działania zabezpieczeń}
\chapter{Ocena przydatności dydaktycznej stanowiska}
\label{chap: Ocena przydatności dydaktycznej stanowiska}

\section{Możliwe ćwiczenia laboratoryjne}
\label{sec: Możliwe ćwiczenia laboratoryjne}

\section{Możliwości rozbudowy stanowiska}
\label{sec: Możliwości rozbudowy stanowiska}

\include{Text/7_Summary}


\bibliographystyle{Config/IEEEtranDOI}
\bibliography{Config/Bibliography}
\vspace*{\fill}
\begingroup
\noindent
Oświadczenie autora projektu:\vspace{15pt}\newline
\noindent
Oświadczam, świadomy odpowiedzialności karnej za poświadczenie nieprawdy, że niniejszy projekt wykonałem osobiście i samodzielnie oraz, że nie korzystałem ze źródeł innych niż wymienione w pracy.\vspace{25pt}\newline
...............................................................................

\hspace{50pt}(Podpis autora pracy)
\endgroup
\vspace*{\fill}





\end{document}
%---------------------------------------------------------------------------


